%% ma: L12-B08-0110
%% lop: 12
%% chuong: 2
%% bai: 8
%% dang: 12.8.27
%% loai: TN4
%% muc: NB
%% dap_an: "C"
%% nguon: "CĐ2.3.pdf giáo viên gửi 27/10/2026 (M07, Lê Minh Khôi – Nguyễn Anh Khôi, Chương II Chủ đề 3 Biểu thức tọa độ của các phép toán vectơ), Câu 21"
%% kien_thuc: "Bài 8 (tọa độ trọng tâm tứ diện)"
%% trang_thai: chinh-thuc
%% ghi_chu: "Đề gốc viết 'trọng tâm của hình chóp SABC'; đã đổi thành 'trọng tâm của tứ diện SABC' cho đúng thuật ngữ SGK"
\begin{ex}
Trong không gian $Oxyz$, cho tứ diện $SABC$ có $S(-1;0;1)$, $A(2;4;-3)$, $B(1;2;-1)$, $C(3;1;-2)$. Tọa độ trọng tâm $G$ của tứ diện $SABC$ là
\choice{$G\left(2;\dfrac72;-2\right)$}{$G\left(\dfrac54;\dfrac74;\dfrac54\right)$}{\True $G\left(\dfrac54;\dfrac74;-\dfrac54\right)$}{$G\left(-2;\dfrac72;-2\right)$}
\loigiai{$G\left(\dfrac{-1+2+1+3}4;\dfrac{0+4+2+1}4;\dfrac{1-3-1-2}4\right)=\left(\dfrac54;\dfrac74;-\dfrac54\right)$.}
\end{ex}
